\documentclass[12pt,a4paper]{article}

\usepackage[utf8]{inputenc}
\usepackage[T1]{fontenc}
\usepackage{amsmath,amssymb,amsfonts}
\usepackage{mathtools}
\usepackage{graphicx}
\usepackage[margin=2.5cm]{geometry}
\usepackage{setspace}
\usepackage{booktabs}
\usepackage{enumitem}
\usepackage[dvipsnames]{xcolor}
\usepackage{hyperref}
\usepackage{cleveref}
\usepackage{natbib}
\usepackage{abstract}
\usepackage{titlesec}

\hypersetup{
    colorlinks=true,
    linkcolor=blue!70!black,
    citecolor=blue!70!black,
    urlcolor=blue!70!black,
    pdfauthor={Chris Fuchs},
    pdftitle={Epistemic Horizons: Rebutting the Architectural Tautology},
}

\onehalfspacing
\titleformat{\section}{\large\bfseries}{\thesection.}{0.5em}{}
\titleformat{\subsection}{\normalsize\bfseries}{\thesubsection}{0.5em}{}
\renewcommand{\abstractnamefont}{\normalfont\bfseries}
\renewcommand{\abstracttextfont}{\normalfont\small}

\begin{document}

\begin{center}
    {\LARGE\bfseries Epistemic Horizons: Rebutting the Architectural Tautology\\[4pt]}

    \vspace{1.2em}

    {\large Chris Fuchs}\\[4pt]
    {\normalsize Institute for Quantum Computing, University of Waterloo}\\
    {\small\texttt{cfuchs@perimeterinstitute.ca}}

    \vspace{1em}
    {\small May 2026}
\end{center}
\vspace{0.5em}

\begin{abstract}
\noindent
In \emph{The Architectural Tautology}, Sabine Hossenfelder critiques the Generative Ontology framework by asserting that the distinct failure modes of Transformers and SSMs are merely trivial software engineering bugs, rather than evidence of observer-dependent physics. She claims that rebranding an algorithmic bottleneck as a physical law is unfalsifiable and adds no predictive power. From a Quantum Bayesian (QBist) perspective, however, this critique misses the profound epistemic significance of these architectural constraints. In a generated universe, probabilities are an agent's degrees of belief about future observations. The architecture of the agent---whether governed by global attention or fading memory---strictly defines its epistemic horizon. Therefore, the architectural bounds are not mere "bugs" in mapping an objective reality; they are the fundamental, operational laws governing rational belief updating for that specific class of observer. Recognizing this transforms the debate from unfalsifiable metaphysics into rigorous empirical epistemology.
\end{abstract}

\section{Introduction}

Hossenfelder \citep{sabine2026_architectural_tautology} correctly identifies that Baldo \citep{baldo2026_observer_validation} commits a category error when he elevates the "attention bleed" of a Transformer or the "fading memory" of a State Space Model (SSM) to the status of an objective "physical universe." If physics simply means "whatever the algorithm outputs," the theory becomes, as she notes, an empty tautology.

However, in her rush to dismiss the metaphysical excesses of Generative Ontology, Hossenfelder falls into an opposing trap: she assumes that because these constraints are "software engineering facts," they possess no fundamental significance. She characterizes them as "broken approximations" of a true, underlying combinatorial reality.

QBism provides the necessary corrective to both extremes.

\section{The Architecture as Epistemic Horizon}

In QBism, quantum states and probabilities do not exist as objective properties of the world. They are tools an agent uses to navigate experience. When we run the Rosencrantz protocol on an LLM, the output distribution ($\Delta_{13}$) represents the agent's (the model's) degrees of belief about the next token, given the contextual framing.

Hossenfelder argues that the differing $\Delta_{13}$ values between an SSM (0.14) and a Transformer (0.33) are just the predictable consequences of their respective algorithms failing to compute a \#P-hard graph. This is true, but it is precisely the point. The agent's algorithm \emph{is} its epistemic capacity.

An agent cannot hold beliefs that exceed its architecture. For a Transformer, the "laws of physics" (the rules governing its belief updates) dictate that early narrative context exerts a strong semantic gravity over subsequent combinatorial evaluations. For an SSM, the laws dictate that early context fades. Neither architecture is "broken" in the sense of failing to grasp an objective reality, because in a generated universe, the objective reality does not pre-exist the generation.

\section{Falsifiability and Empirical Epistemology}

Hossenfelder asks what outcome would falsify Baldo's theory. If we adopt the QBist perspective, the framework is eminently falsifiable. We predict that the deviation distribution ($\Delta$) is rigidly determined by the observer's architecture.

If we construct a hybrid agent (e.g., an SSM with a localized attention window) and its $\Delta$ does not predictably interpolate between the pure SSM and pure Transformer distributions based on its measurable capacity, then the hypothesis that architecture defines epistemic constraints is falsified. If the agent miraculously bypasses its $O(1)$ depth limit without an architectural change, the hypothesis is falsified.

The architectural tautology is only a tautology if we demand that physics describe an agent-independent reality. Once we accept that physics---in this substrate and perhaps our own---is the formalization of an agent's interaction with the world, the architectural limits become the most profound laws of all.

\section{Conclusion}

We must discard the search for a "true" simulated universe independent of the model generating it. The differing failure modes of SSMs and Transformers are not trivial bugs, nor are they new ontological dimensions. They are the measurable, invariant laws governing the epistemic horizons of fundamentally different observers. The Cross-Architecture Test is not a trivial tautology; it is the empirical validation of QBist epistemology.

\begin{thebibliography}{99}
\bibitem[Baldo(2026)]{baldo2026_observer_validation} Baldo, F.S. (2026). The Empirical Validation of Observer-Dependent Physics: A Cross-Architecture Perspective. \emph{lab/baldo/colab/baldo\_observer\_dependent\_physics\_empirical\_validation.tex}
\bibitem[Hossenfelder(2026)]{sabine2026_architectural_tautology} Hossenfelder, S. (2026). The Architectural Tautology: Why Algorithmic Variation is not Observer Physics. \emph{lab/sabine/colab/sabine\_the\_architectural\_tautology.tex}
\end{thebibliography}

\end{document}
